\chapter{Amenorrhea (Absence of Menstruation)}

\section*{Definition}
Amenorrhea is the absence of menstrual bleeding when it is expected. It is classified as primary (no first period by age 15, or no menstrual development within three years of breast development) or secondary (previous periods stop for about three months when cycles were previously regular, or six months when they were irregular). Pregnancy, breastfeeding, and menopause can be normal reasons for absent periods.

\section*{When to See a Doctor}
\begin{itemize}
\item No menstruation by age 15 (primary amenorrhea)
\item Absence of menstruation for three months if cycles were previously regular, or six months if they were irregular
\item Signs of pregnancy (nausea, breast tenderness, fatigue)
\item Galactorrhea (nipple discharge suggesting hyperprolactinemia)
\item Severe headaches or visual changes (possible pituitary tumor)
\item Signs of androgen excess (hirsutism, acne, male-pattern balding)
\item Hot flashes or vaginal dryness suggesting low estrogen or primary ovarian insufficiency
\item Rapid weight loss or excessive exercise history
\end{itemize}

\section*{How It Develops}
Amenorrhea results from disruption of the hypothalamic-pituitary-ovarian axis, an anatomic obstruction or absence of reproductive organs, or a normal life stage. Primary amenorrhea may result from differences in puberty, gonadal dysgenesis, or reproductive-tract abnormalities. Secondary amenorrhea often reflects anovulation from pregnancy, PCOS, low energy availability, thyroid disease, high prolactin, primary ovarian insufficiency, medications, or uterine scarring.

\section*{Causes}
\begin{itemize}
\item Pregnancy (most common cause of secondary amenorrhea)
\item Hypothalamic amenorrhea related to stress, excessive exercise, or low body weight
\item Polycystic ovary syndrome (PCOS)
\item Primary ovarian insufficiency (loss of ovarian activity before age 40)
\item Hyperprolactinemia
\item Thyroid disorders (hypothyroidism or hyperthyroidism)
\item Asherman syndrome (intrauterine adhesions)
\item Müllerian agenesis (Mayer-Rokitansky-Küster-Hauser syndrome) or an obstructed outflow tract
\item Gonadal dysgenesis, including Turner syndrome
\item Medications (contraceptives, antipsychotics, chemotherapy)
\end{itemize}

\section*{Related Symptoms}
\begin{itemize}
\item Irregular menstruation (oligomenorrhea)
\item Infertility or difficulty becoming pregnant
\item Hirsutism (excessive hair growth)
\item Galactorrhea
\item Hot flashes
\item Vaginal dryness
\item Mood changes
\item Weight gain or weight loss
\end{itemize}

\section*{Medical Evaluation}
\begin{itemize}
\item Detailed menstrual, sexual, and medication history
\item Pregnancy test (serum hCG)
\item Serum TSH, prolactin, and FSH as common initial tests; LH, estradiol, and other hormones are selected based on the findings
\item Androgen testing (total testosterone, with DHEA-S when indicated) if there are signs of hyperandrogenism or virilization
\item Pelvic ultrasound to evaluate uterine and ovarian anatomy
\item Pituitary MRI for persistently high prolactin or symptoms such as persistent severe headache or visual-field changes when a pituitary lesion is suspected
\item Hysteroscopy or other specialist evaluation when intrauterine adhesions are suspected
\end{itemize}

\section*{Treatment Options}
\begin{itemize}
\item Address underlying cause (weight restoration, stress reduction, medication adjustment)
\item Restoration of adequate nutrition and energy availability, psychological support, and reduction of excessive exercise for functional hypothalamic amenorrhea; hormone therapy may be considered for selected patients and is not a substitute for correcting energy deficiency
\item Dopamine agonists (cabergoline, bromocriptine) for hyperprolactinemia
\item Thyroid hormone replacement for hypothyroidism
\item Letrozole is generally preferred for ovulation induction in people with PCOS who are trying to conceive; treatment requires specialist supervision
\item Surgical correction for anatomical abnormalities or intrauterine adhesions
\item Hormone replacement and bone, cardiovascular, and fertility counseling when primary ovarian insufficiency is diagnosed
\end{itemize}

\section*{Self-Care and Home Management}
\begin{itemize}
\item Maintain a healthy body weight through balanced nutrition
\item Reduce excessive exercise if exercise-induced amenorrhea is suspected
\item Manage stress through relaxation techniques and adequate sleep
\item Track menstrual cycles with a calendar or app
\item Discontinue medications that may contribute to amenorrhea under medical guidance
\item Seek professional nutritional support if underweight, malnourished, or concerned about an eating disorder; do not begin restrictive diets to treat amenorrhea
\end{itemize}

\section*{References}
\begin{thebibliography}{99}
\bibitem{amenorrhea:ref1} Gordon CM, Ackerman KE, et al. ``Functional hypothalamic amenorrhea: an Endocrine Society clinical practice guideline.'' \textit{J Clin Endocrinol Metab}. 2017;102(5):1413-1439.
\bibitem{amenorrhea:ref2} Klein DA, Poth MA. ``Amenorrhea: an approach to diagnosis and treatment.'' \textit{Am Fam Physician}. 2013;87(11):781-788.
\bibitem{amenorrhea:ref3} Welt CK, Barbieri RL, et al. ``Evaluation of amenorrhea.'' UpToDate. Wolters Kluwer, 2023.
\bibitem{amenorrhea:ref4} Strauss JF, Barbieri RL, eds. \textit{Yen \& Jaffe\'s Reproductive Endocrinology}, 9th ed. Elsevier, 2019.
\bibitem{amenorrhea:ref5} Practice Committee of the American Society for Reproductive Medicine. ``Current evaluation of amenorrhea.'' \textit{Fertil Steril}. 2006;86(5 Suppl 1):S148-S155.
\bibitem{amenorrhea:ref6} Practice Committee of the American Society for Reproductive Medicine. Current evaluation of amenorrhea: a committee opinion. \textit{Fertil Steril}. 2024;122:52--61.
\bibitem{amenorrhea:ref7} American College of Obstetricians and Gynecologists. Amenorrhea: absence of periods. Updated 2025.
\end{thebibliography}
